% AUTO-GENERATED by scripts/ulm7883227_pala_disjoint_support_table.py; DO NOT EDIT BY HAND.
\begin{table}[t]
\centering
\small
\resizebox{\linewidth}{!}{%
\begin{tabular}{l c c c}
\hline
Score & AUC (H1 vs H0) & TPR@$10^{-3}$ & shell FPR @70\% core recall \\
\hline
Baseline (power Doppler) & 0.502 [0.495,0.508] & 2.6e-4 [5.9e-5,5.5e-4] & 0.704 [0.695,0.713] \\
Baseline (Kasai lag-1 magnitude) & 0.500 [0.494,0.507] & 2.0e-5 [0,5.9e-5] & 0.709 [0.701,0.718] \\
Fixed matched-subspace head & 0.510 [0.503,0.517] & 3.7e-4 [1.8e-4,6.1e-4] & 0.709 [0.700,0.719] \\
Adaptive head & 0.512 [0.505,0.519] & 3.7e-4 [1.8e-4,6.1e-4] & 0.704 [0.693,0.716] \\
Fully whitened specialist & 0.518 [0.511,0.525] & 0.005 [0.003,0.006] & 0.695 [0.684,0.707] \\
\hline
\end{tabular}%
}
\caption{Disjoint-source robustness variant of the ULM localization-derived structural audit. Localization-derived reference construction uses odd source blocks only, while scored windows come from even source blocks only, under the same frozen 64-frame no-registration profile used in the main audit. Entries report window-level means with 95\% bootstrap intervals over the 70 evaluated windows. At block level, the fully whitened specialist has higher AUC than PD, Kasai, fixed, and adaptive on 10/10, 10/10, 9/10, and 7/10 evaluation blocks, respectively.}
\label{tab:ulm7883227_structural_roc_disjoint_supp}
\end{table}
